% ============================================================
% Chapter 7: Results
%   Section: RQ4 --- The System
% Included from chapters/07-results/results.tex
% Sources: backend/data/evaluation/queue/ (claim selector precision,
% recomputed 2026-10-16: 136 labelled, 78 skipped; bootstrap 10,000,
% seed 2026); the X-Fact reports; docs/decisions/deployment.md and
% selection.md (time per article and per round).
% ============================================================

\section{RQ4 --- How Much of the Work the System Takes On}
\label{sec:results-rq4}

RQ4 asks, of the three steps joined, what share of a journalist's work
the system takes on and at what quality. Its answer is a profile, built
step by step from the other questions' measurements and from three of
its own.

\paragraph{Choosing the claims.} Over the nine daily labelling batches
from 29~September to 14~October, the claim selector proposed 214
claims, of which 136 were labelled and 78 rejected: a precision of
0.636 [0.570, 0.701] (bootstrap over claims). It is markedly higher in
English (98 of 141, 0.70) than in Spanish (38 of 73, 0.52). The
rejection reasons were opinion or interpretation (22), a fragment of a
sentence (21), a trivial statement (15), a prediction (8), a duplicate
(6) and others (6). Two cautions apply. The batches from early October
onwards were worked by the AI labelling pass, so the figure mixes the
annotator's decisions with the AI's and will be recomputed on the
reviewed batches; and the claim splitter was changed on 9~October,
after fragments cut at abbreviations were found, so the later batches
measure a slightly different selector. About two in three proposed
claims are worth checking; the system takes on the choice of claims,
but a person still discards about one in three.

\paragraph{Verifying them.} On the X-Fact pilot the verifier gave a
definitive verdict for 95.3\% of claims (its coverage) and was right on
31.1\% of those (its selective accuracy): it almost never abstains, and
when it answers it is mostly wrong. On that evidence the 3B verifier
cannot take over verification on political claims; whether it can on
constructive news, and what share of its verdicts would be usable
without re-checking, is the custom-set measurement of RQ2.
\pending{coverage, selective accuracy, the share usable without
re-checking and the confidently-wrong share on the custom test split.}

\paragraph{Time.} A claim took 36.5~s at the median on the development
laptop in the pilot (90th percentile 59.9~s), of which 14.8~s
retrieval, 8.6~s ranking and a 5.6~s model call; on the CPU-only
production stack a four-claim article took 184~s, with the model 180 of
those seconds. A round of discovery takes about a minute; the AI
selection about 70~s for 117 candidates with the 3B model. Against
these, the time a person spends is not yet measured: the labelling
tool does not record time per fact.
\pending{time per fact from the labelling tool (median and
interquartile range), against the system's time per claim; the effort
estimate of Subsection~\ref{subsec:harness-effort}.}

\paragraph{The profile so far.} Finding: the system reads a tenth to
two fifths of the links and orders them so that the first twenty
places hold few stories not worth offering, on AI labels
(Section~\ref{sec:results-rq1}). Choosing claims: about two in three of
its proposals are checkable. Verifying: not yet at a quality a
journalist could rely on with the 3B model, on the only data measured.
At every step a person is still needed for the judgement the system is
weakest at --- hard news against constructive news on the same topic,
opinion against fact in Spanish, and whether a verdict can be trusted.
